\documentclass[a4paper]{ltjsarticle}
\usepackage{./files/newwordtest}
\usepackage[margin=15mm]{geometry}
\usepackage{luatexja-ruby}
\renewcommand{\anscolor}{red}
\pagestyle{empty}
\begin{document}
\leftskip=1em
\TitleHead{3}

\section{\blank に当てはまる適切な単語または表現を選択肢の中から一つ選びなさい。}
\nJtoE{579}{1}{shelter}{cellar}{felter}{welter}{1}
\nJtoE{531}{1}{publicity}{puberty}{pungency}{pulsation}{1}
\nJtoE[-1]{411}{4}{consort}{content}{consult}{consent}{4}
\nJtoE[-1]{575}{2}{shaved}{shaded}{shed}{shattered}{2}
\nJtoE[-1]{518}{3}{subtle}{submissive}{substantial}{suburban}{3}

\section{下線部の単語または表現を和訳しなさい。}
\nEtoJ{560}{}{F}{}% 
\nEtoJ{404}{先行する}{F}{}% 
\nEtoJ{508}{}{F}{}% 
\nEtoJ{589}{}{F}{}% 
\nEtoJ{405}{実行する}{F}{}% 
\nEtoJ{520}{広い}{T}{}% 
\nEtoJ{431}{実行する}{T}{}% 
\nEtoJ{572}{えつらんする}{T}{日本語的に「見て回る」など}% 
\nEtoJ{514}{レシート}{T}{}% 
\nEtoJ{529}{致命的な}{F}{そういう意味もある}% 
\nEtoJ{567}{柱}{F}{}% 
\nEtoJ{551}{なびかない}{F}{}% 
\nEtoJ{416}{従事}{F}{}% 
\nEtoJ{539}{}{F}{}% 
\nEtoJ{437}{圧}{F}{}% 
\nEtoJ{490}{したがらない}{T}{}% 
\nEtoJ{570}{}{F}{}% 
\nEtoJ{401}{解決}{T}{}% 
\nEtoJ{410}{見とおし}{T}{}% 
\nEtoJ[0.9]{462}{計画する}{T}{そういう意味もある}% 

\section{次の各文の\blank に最もよく当てはまる語または表現を答えなさい。}
\nEW{412}{}{F}{}% 
\nEW[-1]{415}{encourage}{F}{}% 
\nEW[-1]{561}{stress}{T}{想定してたのはstrain}% 
\nEW[-1]{475}{}{F}{}% 
\nEW[-1]{550}{}{F}{}% 
\addtocounter{correct}{-1}
\refstepcounter{correct}{\label{count:score}}
\end{document}